% page 335 of the facsimile (image p-334 of the scan)
% block 160.0 x 254.9 mm ; left margin 42.0 mm
\pgc{20.320mm}{0.000mm}{
\l{}
\l{}
\l{}
\l{\cel{50}327}
\l{}
\l{lacerar\cel{2}(\sou{trans.)}\cel{3}Dispecigar arache, tale ke lu ne plus}
\l{\cel{3}esos uzebla. - EFIS.}
\l{}
\l{lacerto.\cel{2}(\sou{zool.)}\cel{4}Reptero sauria, kun quar gambi kurta e}
\l{\cel{3}tenua, kaudo longa ek ringi flexebla e facile ruptebla (ma}
\l{\cel{3}ri-kreskas). L. lacerta. - EFIS.}
\l{}
\l{laciva. Inklinata a la plezuri amorala ; qua su abandonas}
\l{\cel{3}grosiere a la plezuri amorala, koitala. - DEFIS.}
\l{}
\l{lado.\cel{2}Fero-tolo dina, kovrita ek induturo de stano, qua pro-}
\l{\cel{3}tektas lu de rusto. - IS.}
\l{}
\l{\sou{ladano}.\cel{2}(\sou{Farmacio}). Gumo rezina aromata, quan donas diversa}
\l{\cel{3}arbusti, precipue la cisto di Kretia. - L. ladanum. - DEFIRS}
\l{}
\l{\sou{lago}.\cel{2}Mareto ek aquo stagnanta, interne di kontinento.- EFIS.}
\l{}
\l{\sou{laguno}. (\sou{geogr,)}\cel{3}Parto di maro, poke profunda, inter insu-}
\l{\cel{3}leti. - DEFIRS.}
\l{}
\l{\sou{laika}.\cel{2}Qua ne resortisas ad eklezio. - DEFIS.}
\l{}
\l{\sou{lako}.\cel{2}I. Materio rezinoza, obtenata de ula arbori, mi-diafana,}
\l{\cel{3}brune-reda, uzata por facar ula vaxi, tintivi, bela vernisi,}
\l{\cel{3}e c. - II. Verniso ek Chinia, tre prizata, nigra o reda,}
\l{\cel{3}ornita ye figuri, arabeski, orizuri. - DEFIRS.}
\l{}
\l{\sou{lakeo}. Servisto livreoza, employata ye akompanar seque sua}
\l{\cel{3}mastro. - DEFIRS.}
\l{}
\l{\sou{lakonika}. Qua expresas lo pensata per tre poka vorti (segun}
\l{\cel{3}la maniero di la habitanti Lakonia, lor la epoki antiqua.}
\l{\cel{3}- DEFIRS.\cel{17}di}
\l{}
\l{lakrimo.\cel{3}Guto de humoro limpida, quan sekrecas glando di la}
\l{\cel{3}okulo, efike da emoceso kordiala o da efekto kozala. - EFIS.}
\l{}
\l{lakto.\cel{2}Liquido opaka, blanka, sukroza, sekrecata da la mamo-}
\l{\cel{3}glandi di la muliero e di la mamiferini, por la nutro di la}
\l{\cel{3}jus-naskinti. - EFIS.}
\l{}
\l{laktazo. (\sou{kemio}). Diastazo, o fermento dissolvebla, sekrecata}
\l{\cel{3}da kelka hefi, e qua transformas la laktoso a glikoso o}
\l{\cel{3}galaktoso. - DEFIRS.}
\l{}
\l{laktometro. Areometro uzata por mesurer la denseso-grado di}
\l{\cel{3}lakto. - DEFIRS.}
\l{\cel{25}12\cel{2}22\cel{2}11}
\l{laktoso. (\sou{kemio}). Sukro C\cel{3}H\cel{3}O\cel{2}, en la lakto di la mami-}
\l{\cel{3}feri. - DEFIRS.}
\l{}
\l{lakuno. I. Spaco vakua en la kontinueso di korpo. - II. Vakuajo}
\l{\cel{3}en texto, en kateniguro, en serio. - EFIS.}
}
